% hopfield_report.tex
\documentclass[11pt,a4paper]{article}
\usepackage[utf8]{inputenc}
\usepackage[T1]{fontenc}
\usepackage{geometry}
\geometry{margin=2.5cm}
\usepackage{amsmath,amssymb}
\usepackage{graphicx}
\usepackage{caption}
\usepackage{listings}
\usepackage{booktabs}
\usepackage{float}
\usepackage{enumitem}

\lstdefinelanguage{Matlab}{
  keywords={break,case,catch,continue,else,elseif,end,for,function,global,if,otherwise,persistent,return,switch,try,while},
  sensitive=true,
  comment=[l]{\%},
  morestring=[b]',
}
\lstset{
  language=Matlab,
  basicstyle=\ttfamily\small,
  numbers=left,
  numberstyle=\tiny,
  stepnumber=1,
  numbersep=6pt,
  frame=single,
  breaklines=true,
  showstringspaces=false,
  captionpos=b
}

\title{Intro AI — HW2}
\author{MohammadAmin HorriFrahani - S.ID = 40216673}
\date{\today}

\begin{document}
\maketitle

\begin{abstract}
Hopfield-network exercise
\end{abstract}

\tableofcontents
\newpage

\section{Problem statement}
Binary Hopfield network with \(N=4\). Store two patterns via Hebbian learning:
\[
p^{(1)}=(+1,+1,-1,-1),\qquad p^{(2)}=(+1,-1,+1,-1).
\]
Weights:
\[
W=\frac{1}{N}\sum_{\mu=1}^2 p^{(\mu)} p^{(\mu)\top},\qquad W_{ii}=0,\quad W=W^\top.
\]
Energy:
\[
E(s)=-\tfrac12 s^\top W s.
\]
Asynchronous update order: \(1\to2\to3\to4\), with
\[
s_i\leftarrow \operatorname{sgn}\big((W s)_i\big),
\]
and \(\operatorname{sgn}(0)\) leaves \(s_i\) unchanged.

Tasks: (a) compute \(W\) explicitly. (b) verify stored patterns are fixed points. (c) from initial \(s^{(0)}=(+1,-1,-1,-1)\) run asynchronous updates for two full sweeps; list state after each single-neuron update and energy after each sweep. (d) starting from \(x=(+1,+1,+1,-1)\) run one sweep and say which prototype basin it enters. (e) short comparison to K-means Voronoi cells.

\section{Solution}

\subsection{(a) Weights}
Compute outer-products and average (then zero the diagonal):
\[
W=\frac{1}{4}\big(p^{(1)}p^{(1)\top}+p^{(2)}p^{(2)\top}\big)
=\begin{pmatrix}
\frac12 & 0 & 0 & -\frac12\\[4pt]
0 & \frac12 & -\frac12 & 0\\[4pt]
0 & -\frac12 & \frac12 & 0\\[4pt]
-\frac12 & 0 & 0 & \frac12
\end{pmatrix}.
\]
Enforce \(W_{ii}=0\):
\[
\boxed{W=\begin{pmatrix}
0 & 0 & 0 & -\tfrac12\\[4pt]
0 & 0 & -\tfrac12 & 0\\[4pt]
0 & -\tfrac12 & 0 & 0\\[4pt]
-\tfrac12 & 0 & 0 & 0
\end{pmatrix}.}
\]

\subsection{(b) Fixed points}
Compute \(Wp^{(1)}\) and \(Wp^{(2)}\):
\[
W p^{(1)}=\begin{pmatrix}0.5\\[2pt]0.5\\[2pt]-0.5\\[2pt]-0.5\end{pmatrix},\quad
W p^{(2)}=\begin{pmatrix}0.5\\[2pt]-0.5\\[2pt]0.5\\[2pt]-0.5\end{pmatrix}.
\]
Applying \(\operatorname{sgn}(\cdot)\) (with \(\operatorname{sgn}(0)\) leaving bit unchanged) yields exactly \(p^{(1)}\) and \(p^{(2)}\). So both stored patterns are fixed points.

\subsection{(c) Recall steps and energy}
Energy formula:
\[
E(s)=-\tfrac12 s^\top W s.
\]

Start \(s^{(0)}=(+1,-1,-1,-1)\). Update order \(1\to2\to3\to4\). I list state \emph{after each single-neuron update}.

\begin{enumerate}[label=\textbf{Update \arabic*:},leftmargin=*,nosep]
  \item After updating neuron 1: \(s=(+1,-1,-1,-1)\).
  \item After updating neuron 2: \(s=(+1,+1,-1,-1)\).
  \item After updating neuron 3: \(s=(+1,+1,-1,-1)\).
  \item After updating neuron 4: \(s=(+1,+1,-1,-1)\).
\end{enumerate}

Energy after sweep 1 (after 4 updates):
\[
E=-1.0.
\]

Second sweep (same order) leaves the state unchanged:
\begin{enumerate}[label=\textbf{Update \arabic*:},start=5,leftmargin=*,nosep]
  \item After updating neuron 1 (sweep 2): \(s=(+1,+1,-1,-1)\).
  \item After updating neuron 2: \(s=(+1,+1,-1,-1)\).
  \item After updating neuron 3: \(s=(+1,+1,-1,-1)\).
  \item After updating neuron 4: \(s=(+1,+1,-1,-1)\).
\end{enumerate}

Energy after sweep 2:
\[
E=-1.0.
\]

Thus after the first sweep the network reached the stored pattern \(p^{(1)}\) and remained there.

\subsection{(d) Prototype link}
Start \(x=(+1,+1,+1,-1)\) and run one asynchronous sweep \(1\to2\to3\to4\):
\begin{itemize}
  \item Updating neuron 1: no change.
  \item Updating neuron 2: its field is negative so \(s_2\) flips to \(-1\), resulting in \(s=(+1,-1,+1,-1)=p^{(2)}\).
  \item Remaining updates keep it at \(p^{(2)}\).
\end{itemize}
So \(x\) enters the basin of prototype \(p^{(2)}\).

\subsection{(e) Conclusion}
Hopfield attraction basins are energy basins shaped by learned correlations; they are not identical to K-means Voronoi cells (which are Euclidean-distance partitions). If prototypes overlap or are highly correlated, Hopfield networks can produce spurious attractors (mixed states) or reduced basins of attraction, harming recall.

\appendix
\section{Appendix A: MATLAB code reproducing computations}
\label{app:matlab}

\begin{lstlisting}[language=Matlab,caption={MATLAB script: compute W, verify fixed points, simulate asynchronous updates}]
% hopfield_sim.m
% Compute weights, verify fixed points, simulate asynchronous updates.

p1 = [1; 1; -1; -1];
p2 = [1; -1; 1; -1];
N = 4;

W = (p1*p1' + p2*p2')/N;
W(1:5:end) = 0; % set diagonal to zero

disp('Weight matrix W:')
disp(W)

% Verify fixed points
wp1 = W * p1;
wp2 = W * p2;
disp('W*p1 ='); disp(wp1')
disp('sign(W*p1) ='); disp(sign(wp1)')
disp('W*p2 ='); disp(wp2')
disp('sign(W*p2) ='); disp(sign(wp2)')

% Energy function
energy = @(s) -0.5 * s' * W * s;

% (c) asynchronous recall from s0
s = [1; -1; -1; -1];
order = [1 2 3 4];
states = zeros(4,8);
k = 1;
for sweep = 1:2
  for idx = order
    h = W(idx,:) * s;
    if h > 0, s(idx) = 1;
    elseif h < 0, s(idx) = -1;
    end
    states(:,k) = s;
    k = k + 1;
  end
  fprintf('After sweep %d, state = [%s], energy = %.2f\n', sweep, num2str(s.'), energy(s));
end
disp('States after each single update (columns):'); disp(states)

% (d) test unseen x
x = [1;1;1;-1];
s = x;
for idx = order
  h = W(idx,:) * s;
  if h > 0, s(idx) = 1;
  elseif h < 0, s(idx) = -1;
  end
end
fprintf('After one sweep from x, state = [%s]\n', num2str(s.'));
\end{lstlisting}

\section{Appendix B: numeric checks}
Energy at the final attractor \(p^{(1)}\): \(E=-1.0\) (numeric).

\subsection*{Question 2: SOM}

The output weights are
\[
\begin{array}{c|rrrrrr}
\text{Node} & A & B & C & D & E & F\\\hline
w_1 & -1 & 0 & 3 & -2 & 3 & 4\\
w_2 & 2 & 4 & -2 & -3 & 2 & -1
\end{array}
\]
Input: \(x=(2,-4)\). Distances (squared) and steps:

\[
\begin{aligned}
d_A^2 &= (2-(-1))^2 + (-4-2)^2 = 3^2 + (-6)^2=9+36=45,\\
d_B^2 &= (2-0)^2 + (-4-4)^2 = 2^2 + (-8)^2=4+64=68,\\
d_C^2 &= (2-3)^2 + (-4-(-2))^2 = (-1)^2 + (-2)^2=1+4=5,\\
d_D^2 &= (2-(-2))^2 + (-4-(-3))^2 = 4^2 + (-1)^2=16+1=17,\\
d_E^2 &= (2-3)^2 + (-4-2)^2 = (-1)^2 + (-6)^2=1+36=37,\\
d_F^2 &= (2-4)^2 + (-4-(-1))^2 = (-2)^2 + (-3)^2=4+9=13.
\end{aligned}
\]

Minimum squared distance is \(d_C^2=5\Rightarrow\) \textbf{winner = C}.

Adaptation rule: \(w' = w + \alpha h (x-w)\) with \(\alpha=0.5\). Neighbourhood \(h=1\) for winner, \(h=0.5\) for immediate neighbours (B and D), \(h=0\) otherwise.

Compute updates:
\[
\begin{aligned}
w_C' &= (3,-2) + 0.5\cdot 1\cdot((2,-4)-(3,-2))=(3,-2)+0.5\cdot(-1,-2)=(2.5,-3.0),\\
w_B' &= (0,4) + 0.5\cdot 0.5\cdot((2,-4)-(0,4))=(0,4)+0.25\cdot(2,-8)=(0.5,2.0),\\
w_D' &= (-2,-3) + 0.25\cdot((2,-4)-(-2,-3))=(-2,-3)+0.25\cdot(4,-1)=(-1.0,-3.25).
\end{aligned}
\]

Nodes \(A,E,F\) unchanged. Thus updated weights:
\[
\begin{array}{c|rrrrrr}
\text{Node} & A & B & C & D & E & F\\\hline
w'_1 & -1 & 0.5 & 2.5 & -1 & 3 & 4\\
w'_2 & 2 & 2.0 & -3.0 & -3.25 & 2 & -1
\end{array}
\]

\paragraph{Comparison (brief).} Feed-forward networks are usually trained supervised (minimize global loss via backprop); SOMs are unsupervised, use competitive learning and neighbourhood updates, and produce topology-preserving low-dimensional maps for clustering/visualization.


\begin{lstlisting}[language=Matlab,caption={MATLAB script}]


% som_update.m
% Determine winner for x and perform neighbourhood weight updates.

% initial weights
W = [ -1,  0,  3, -2,  3,  4 ;   % row 1 = w1
       2,  4, -2, -3,  2, -1 ]; % row 2 = w2

x = [2; -4];

% compute squared distances
d2 = zeros(1,6);
for j=1:6
    d2(j) = (x(1)-W(1,j))^2 + (x(2)-W(2,j))^2;
end
disp('Squared distances to nodes A..F:')
disp(d2)

% find winner
[~, winner] = min(d2); % index 1..6 corresponds to A..F
nodes = 'ABCDEF';
fprintf('Winner: node %c\n', nodes(winner));

% neighbourhood: ring A-B-C-D-E-F-A
% immediate neighbours of C (index 3) are 2 and 4
alpha = 0.5;
h = zeros(1,6);
h(winner) = 1;
% set immediate neighbours (modular ring)
left = mod(winner-2,6)+1;
right = mod(winner,6)+1;
h(left) = 0.5;
h(right) = 0.5;

% apply updates
Wnew = W;
for j=1:6
    if h(j) ~= 0
        Wnew(:,j) = W(:,j) + alpha * h(j) * (x - W(:,j));
    end
end

disp('Updated weights (columns A..F):')
disp(Wnew)


\end{lstlisting}



\end{document}
